Moreover, from Lemma \ref{lem:baileylattices}, we deduce a very general theorem transforming bilateral Bailey pairs relative to $a$ into bilateral Bailey pairs relative to $aq^{-N}$. Recall the $M$-th elementary symmetric polynomial in $N$ variables defined for $0 \leq M \leq N$ as
\begin{equation*}e_M(X_1, \dots , X_N)= \sum_{1\leq i_1<i_2<\dots<i_M\leq N} X_{i_1} X_{i_2}\cdots X_{i_M},
\end{equation*}
and $e_M(X_1, \dots , X_N)=0$ if $M<0$ or $M>N$.

Recall also that for all $0\leq j \leq N$, the $q$-binomial coefficient is defined by
\begin{equation*}\left[{N\atop j}\right]_q=\left[{N\atop j}\right]:=\frac{(q)_N}{(q)_j(q)_{N-j}}.
\end{equation*}
We extend this definition to $j<0$ and $j>N$ by setting $\bigl[{N\atop j}\bigr]=0$, which, as $N\geq0$, is consistent with the definition of $q$-Pochhammer symbols with negative indices given above.
The general theorem can be stated as follows.

\begin{Theorem}[new bilateral Bailey lattice in higher dimension]\label{thm:multibaileylattice}
Let $(\alpha_n, \beta_n)$ be a bilateral Bailey pair relative to $a$. For all $N\geq 1$, define the pair \smash{$\bigl(\alpha^{(N)}_n, \beta^{(N)}_n\bigr)$} by
\begin{equation}
\alpha^{(N)}_n=\frac{\bigl(1-aq^{2n-N}\bigr)\bigl(aq^{1-N}\bigr)_N}{(1-b_1)\cdots(1-b_N)}\sum_{j \in \Z}(-1)^j\frac{q^{jn-j(j+1)/2}f_{N,j,n}(b_1,\ldots,b_N)}{\bigl(aq^{2n-N-j}\bigr)_{N+1}} \alpha_{n-j},\label{eq:multibaileylatticealpha}
\end{equation}
where
\begin{align}
f_{N,j,n}(b_1,\ldots,b_N) :={}& \sum_{M \in \Z} \sum_{u \in \Z} a^{u}q^{(M-j+u)(n-j+u)+u(n-N)} \begin{bmatrix}
M\\
j-u
\end{bmatrix} \begin{bmatrix}
N-M\\
u
\end{bmatrix}\nonumber\\
&\times (-1)^Me_M(b_1, \dots , b_N),\label{eq:pascalab}
\end{align}
and
\begin{equation}
\beta^{(N)}_n=\left(\prod_{i=1}^N \frac{1-b_iq^n}{1-b_i}\right)\beta_n.\label{eq:multibaileylatticebeta}
\end{equation}
Then, \smash{$\bigl(\alpha^{(N)}_n, \beta^{(N)}_n\bigr)$} is a bilateral Bailey pair relative to $aq^{-N}$.
\end{Theorem}

\begin{Remark}
When $j<0$ or $j > N$, we have $f_{N,j,n}(b_1,\ldots,b_N) =0$ because of the $q$-binomial coefficients in~\eqref{eq:pascalab}. Therefore, the sum in \eqref{eq:multibaileylatticealpha} is actually finite.
\end{Remark}

\begin{Remark}
Note that the sums over $M$ and $u$ in \eqref{eq:pascalab} could equivalently be taken from~$0$ to $N$ and from~$0$ to $j$, respectively. Indeed the $q$-binomial coefficients or $e_M(b_1,\dots ,b_N)$ naturally cancel outside of these ranges. However, the expression with infinite sums makes future calculations easier to write.
\end{Remark}

